\section{VRA Background and Layer Architecture}
\label{sec:background}

\subsection{Section 2 and the Gingles Factors}

Section~2 of the Voting Rights Act of 1965 prohibits electoral practices that
result in the denial or abridgement of the right to vote on account of race,
colour, or membership in a language minority group.
In redistricting, Section~2 is enforced through the three-part test established
in \emph{Thornburg v.\ Gingles}~\citep{gingles1986}:
\begin{enumerate}
  \item The minority group must be sufficiently large and geographically compact
        to constitute a majority in a single-member district.
  \item The minority group must be politically cohesive.
  \item The majority must vote sufficiently as a bloc to enable it to usually
        defeat the minority's preferred candidate.
\end{enumerate}
When all three Gingles preconditions are met, the jurisdiction must draw a
majority-minority district: one in which the minority group constitutes a
majority of the voting-age population (VAP).
The standard threshold is $50\%$ minority VAP, with BVAP (Black VAP) and
HVAP (Hispanic VAP) being the most commonly litigated measures.

\subsection{Majority-Minority District Identification}

Let $G = (V, E)$ be the census-tract adjacency graph for a state with $|V| = n$
tracts.
For each tract $v \in V$, let $\mathrm{pop}(v)$ be the total population and
$\mathrm{mvap}(v)$ be the minority voting-age population (BVAP, HVAP, or a
combined measure, depending on the state).
The \emph{minority VAP fraction} of a district $D \subseteq V$ in plan $\pi$ is:
\[
  \vap(D, \pi) \;=\; \frac{\sum_{v \in D} \mathrm{mvap}(v)}{\sum_{v \in D}
  \mathrm{pop}(v)}.
\]
A district $D$ is a \emph{majority-minority district} under threshold $\theta$
if $\vap(D, \pi) \geq \theta$.
The default threshold in this implementation is $\theta = 0.50$, reflecting
the standard Gingles requirement.

For VRA-aware sampling, we define the \emph{protected set} of a plan $\pi_0$:
\[
  \Prot(\pi_0) \;=\; \bigl\{d \in [k] : \vap(\pi_0^{-1}(d),\, \pi_0) \geq \theta\bigr\}.
\]
This set is computed once at chain construction and does not change as the
chain runs.
The legal justification is that the VRA obligation is defined with respect to
the submitted plan: a majority-minority district in $\pi_0$ must remain a
majority-minority district in every alternative plan produced by the chain.

\subsection{The Two-Layer Architecture}

The redistricting compositor separates algorithm choices into three orthogonal
layers: structure (bisection tree shape), weights (METIS signal), and search
(seed selection strategy).
VRA compliance can be enforced at either the structure layer or the search
layer, with distinct semantics:

\begin{center}
\begin{tabular}{lll}
\toprule
\textbf{Property} & \textbf{VRASection (D.0)} & \textbf{\vr{} (G.13)} \\
\midrule
Layer & Structure & Search \\
Mechanism & METIS edge-weight bias & Hard MH rejection \\
Output & Single plan & Distribution over plans \\
Guarantee & VRA-compliant starting plan & VRA-compliant ensemble \\
Use case & Generate initial plan & Litigation ensemble \\
\bottomrule
\end{tabular}
\end{center}

\noindent
\vraSection{} biases the recursive bisection toward cuts that produce
majority-minority districts; it solves the problem of \emph{finding} a
compliant plan.
\vr{} enforces VRA compliance as an invariant of a Markov chain; it solves
the problem of \emph{sampling} the space of compliant plans.

Neither replaces the other.
A chain that starts from a non-compliant plan will have an initial rejection
burst while it first reaches a compliant state, but since any proposed step
that would move away from compliance is rejected, the chain can never leave
the compliant sub-space once it has entered it.
In practice, using \vraSection{} to generate the starting plan eliminates this
burn-in entirely.

\subsection{VAP Fraction Computation}

The minority VAP data is loaded through the existing vertex-weight pipeline
via \texttt{--weights-override vra-aligned}.
This produces a vector $\mathbf{m} \in [0,1]^n$ of per-tract minority VAP
fractions (i.e., $\mathrm{mvap}(v) / \mathrm{pop}(v)$ for each tract $v$).
The district-level VAP fraction is then a population-weighted average:
\[
  \vap(D, \pi) \;=\;
  \frac{\sum_{v \in D} m_v \cdot \mathrm{pop}(v)}{\sum_{v \in D} \mathrm{pop}(v)}.
\]
The \texttt{minority\_col} parameter selects the minority group column:
\texttt{BVAP} (Black), \texttt{HVAP} (Hispanic), or \texttt{HVAP+BVAP}
(combined, computed at load time).
The values are pre-normalised to $[0,1]$ before being passed to
\texttt{VraRecomChain::new}, so the VAP fraction computation within the
chain reduces to a population-weighted average of scalar values.

The per-step VAP check is $O(|D|)$ per protected district, where $|D|$ is
the number of tracts in the district.
For the states studied (NC, GA, TX), with $|D| \approx n/k$ and
$k \leq 38$, this is negligible compared to the Wilson UST call cost of
$O(|R|^2)$ per step.
\end{antml:parameter>
</invoke>